% q0034.tex — OUR Street
\question{
  id        = {0034},
  title     = {OUR Street},
  source    = {OBECON},
  year      = {2026},
  round     = {Seletiva IEO -- Phase C},
  language  = {English},
  topic     = {Micro},
  subtopic  = {Public Goods},
  type      = {objective},
  statement = {%
    A local community is deciding whether to fund a street lighting system.
    Once installed, the lighting benefits all residents equally. Each resident
    prefers having street lights to not having them, but does not value them
    enough to pay for the entire cost individually. Which of the following best
    explains why a purely voluntary market is likely to underprovide street
    lighting?
  },
  choices   = {%
    \item Because street lighting has increasing marginal costs, private firms
          can't supply it efficiently.
    \item Because individuals don't consider the benefits that street lighting
          provides to others when deciding how much to contribute.
    \item Because street lighting is non-rival but excludable, firms can't
          charge an efficient price.
    \item Because lights are an inferior good and therefore have little
          perceived value to consumers.
  },
  answer    = {B},
  solution  = {%
    Street lighting is a non-rival, non-excludable public good. Individuals
    have incentives to free-ride, ignoring the positive externalities their
    contributions create for others, leading to underprovision. (B) is correct.
  },
}
